\begin{center}
{\LARGE\bfseries From Gaussians to a Cluster Mass Map}\\[6pt]
{\large A worksheet on likelihood stacking, reference shears and $\kappa$ reconstruction\ifanswers\ \textemdash\ \textcolor{ForestGreen!60!black}{Answer key}\fi}\\[4pt]
{\normalsize KL-NN paper, Sec.~6.4--6.5 \quad\textbullet\quad October 2026}
\end{center}

\section*{How to use this worksheet}
\begin{itemize}[itemsep=2pt]
  \item The problems build on each other. Each one has a \textbf{Given} box with everything you need that is not a
    result of an earlier problem. A ``You'll need from before'' line names the earlier results it uses.
  \item Every Given box names its \textbf{source}: a paper, section and equation number, read in the source text. A
    result that we derived ourselves and that is not in those sources is marked \emph{derived here}. The full list of
    sources is at the end, and \texttt{references.md} quotes the line each citation points to.
  \item Problems marked \py\ have a companion cell in \texttt{cluster\_worksheet.ipynb} (numpy, scipy, matplotlib). Do
    the pen-and-paper part first, then check it numerically.
  \item Parts A--B are statistics, C--D are lensing and Fourier analysis, E--F are regularised inversion, and G is the
    cluster analysis itself, assembled from the earlier parts.
\end{itemize}

\subsection*{Notation}
\begin{center}
\begin{tabular}{@{}ll@{}}
\toprule
$\vg = (g_1, g_2)$ & reduced shear of a galaxy; $\gamma$ shear, $\kappa$ convergence \\
$L_i(\vg)$ & shear likelihood of galaxy $i$ (its NPE posterior divided by the training prior) \\
$C_i$, $F_i = C_i^{-1}$ & covariance and precision (inverse covariance) of $L_i$ \\
$R_i$ & rotation from the image frame to galaxy $i$'s tangential/cross frame \\
$\vG$, $\hat\vG$ & common shear of a cell; its stacked estimate (posterior mean) \\
$\gref$ & reference shear of a cell (Problem B4) \\
$m$, $c$ & multiplicative and additive bias, $\hat g = (1+m)g + c$ \\
$\mathsf H$, $\mathsf H_d$, $\mathsf H_p$ & Hessian of the negative log posterior; its data and prior parts \\
$\mathsf R$ & resolution matrix, $\E[\hat\kappa] = \mathsf R\,\kappa_{\rm true}$ \\
\bottomrule
\end{tabular}
\end{center}

\tableofcontents
\clearpage

\section*{Part A. Combining measurements}\addcontentsline{toc}{section}{Part A. Combining measurements}
% ============================== Part A ==============================
\problem{A1}{Two measurements of one number}

\begin{given}
A measurement $d$ of an unknown $x$ with Gaussian noise of standard deviation $\sigma$ has likelihood
\[
L(x) \propto \exp\!\left[-\frac{(x-d)^2}{2\sigma^2}\right].
\]
Independent measurements multiply: the joint likelihood of $d_1$ and $d_2$ is $L_1(x)\,L_2(x)$. Maximising it is the
same as minimising $\chi^2 = \sum_i (d_i - x)^2/\sigma_i^2$, the weighted least-squares problem.
\source{\cHBLwls.}
\end{given}

\need{nothing; this is the start.}

\begin{enumerate}
  \item Show that $L_1(x)L_2(x)$ is again a Gaussian in $x$. Find its mean $\hat x$ and variance $\sigma_{\hat x}^2$.
  \item Write $\hat x$ as a weighted average of $d_1$ and $d_2$. What are the weights?
  \item A KL galaxy measures one shear component as $d_1 = 0.03 \pm 0.05$; a second galaxy gives $d_2 = 0.01 \pm 0.02$.
    Compute $\hat x$ and $\sigma_{\hat x}$. Which galaxy dominates, and why?
\end{enumerate}

\solution{
(a) Adding the exponents,
$-\tfrac12\big[(x-d_1)^2/\sigma_1^2 + (x-d_2)^2/\sigma_2^2\big] = -\tfrac12 (w_1+w_2)\,x^2 + (w_1d_1 + w_2d_2)\,x + \text{const}$,
with $w_i = 1/\sigma_i^2$. Completing the square gives a Gaussian with
\[
\hat x = \frac{w_1 d_1 + w_2 d_2}{w_1 + w_2},\qquad \sigma_{\hat x}^2 = \frac{1}{w_1 + w_2}.
\]
(b) The weights are the precisions $w_i = 1/\sigma_i^2$, normalised to sum to 1. Precisions add.
(c) $w_1 = 400$, $w_2 = 2500$: $\hat x = (12 + 25)/2900 = 0.0128$ and $\sigma_{\hat x} = 2900^{-1/2} = 0.0186$. The second
galaxy dominates: its weight is $(0.05/0.02)^2 = 6.25$ times larger. The combined error is smaller than either.
}

% --------------------------------------------------------------------
\problem{A2}{Many measurements}

\begin{given}
The result of A1 generalises to $N$ independent Gaussian measurements $d_i \pm \sigma_i$ of one $x$: the product of
the likelihoods is Gaussian, with precision $\sum_i w_i$ and mean $\sum_i w_i d_i / \sum_i w_i$, where
$w_i = 1/\sigma_i^2$.
\source{\cHBLwls.}
\end{given}

\need{A1 (precisions add; the mean is precision-weighted).}

\begin{enumerate}
  \item Show that for equal errors $\sigma_i = \sigma$ the combined error is $\sigma/\sqrt N$.
  \item A KL galaxy measures each shear component with $\sigma_g \approx 0.05$. How many such galaxies are needed for
    an error of $0.01$? What is the error for a cell of 5 galaxies?
  \item Our cluster cells hold about 5 galaxies. The subhalo changes the shear in most cells by $\lesssim 0.01$. Can a
    single cell detect the subhalo?
\end{enumerate}

\solution{
(a) $\sum_i w_i = N/\sigma^2$, so $\sigma_{\hat x} = \sigma/\sqrt N$.
(b) $0.05/\sqrt N = 0.01$ gives $N = 25$. With 5 galaxies, $0.05/\sqrt5 = 0.022$. This is the per-cell error
seen in the cluster shear maps.
(c) No: a $\lesssim 0.01$ signal against a $0.022$ error is below $0.5\sigma$ per cell. Detecting the subhalo needs
the information of many cells (or many galaxies) combined. That is what the $\kappa$ map does (Parts E--G).
}

% --------------------------------------------------------------------
\problem{A3}{Two-component measurements with covariance matrices}

\begin{given}
A measurement $\vd_i$ of a 2-vector $\vx$ (for us, $\vx = (g_1, g_2)$) with Gaussian noise of covariance $C_i$ has
likelihood $L_i(\vx) \propto \exp\!\big[-\tfrac12 (\vx-\vd_i)^\T F_i\,(\vx - \vd_i)\big]$, with precision matrix
$F_i = C_i^{-1}$. The general weighted least-squares solution, for data $Y$ with covariance $C$ and a linear model
$Y = AX$, is $\hat X = [A^\T C^{-1} A]^{-1} A^\T C^{-1} Y$, with covariance $[A^\T C^{-1} A]^{-1}$.
\source{\cHBLwls.}
\end{given}

\need{A1, A2.}

\begin{enumerate}
  \item Show that $\prod_i L_i(\vx)$ is Gaussian with precision $\Ftot = \sum_i F_i$ and mean
    $\hat\vx = \Ftot^{-1}\sum_i F_i \vd_i$. (Either complete the square, or use the weighted least-squares solution
    with $A$ a stack of $2\times2$ identity matrices.)
  \item Two galaxies: $\vd_1 = (0.02, 0)$ with $F_1 = \mathrm{diag}(2500, 100)$; $\vd_2 = (0, 0.02)$ with
    $F_2 = \mathrm{diag}(100, 2500)$. Compute $\hat\vx$. Compare with a ``scalar'' average that weights each galaxy by
    the trace of its $F_i$.
  \item Now $\vd_1 = (0.02, 0)$ with $F_1 = \begin{pmatrix}1300 & 1200\\ 1200 & 1300\end{pmatrix}$ and
    $\vd_2 = (0,0)$ with $F_2 = 100\,I$. Compute $\hat\vx$. Neither measurement has a non-zero second component; why
    is $\hat x_2 \neq 0$?
  \item KL kinematics measure one shear direction much better than the other. What does this problem say about
    combining KL galaxies?
\end{enumerate}

\solution{
(a) $\sum_i (\vx - \vd_i)^\T F_i (\vx - \vd_i) = \vx^\T \Ftot \vx - 2\vx^\T \sum_i F_i\vd_i + \text{const}
= (\vx - \hat\vx)^\T \Ftot (\vx - \hat\vx) + \text{const}$, with $\hat\vx = \Ftot^{-1}\sum_i F_i\vd_i$.
(b) $\Ftot = \mathrm{diag}(2600, 2600)$ and $\sum F_i\vd_i = (50, 50)$, so $\hat\vx = (0.0192, 0.0192)$. Each galaxy
dominates the component it measures well. The trace weights are equal (2600 each), giving $(0.01, 0.01)$: half the
signal is lost, because the scalar average throws away which direction each galaxy measures.
(c) $\Ftot = \begin{pmatrix}1400 & 1200\\1200 & 1400\end{pmatrix}$ and $\sum F_i\vd_i = (26, 24)$. Then
$\Ftot^{-1} = \frac{1}{520000}\begin{pmatrix}1400 & -1200\\ -1200 & 1400\end{pmatrix}$, so
$\hat\vx = (0.0146, 0.0046)$. $F_1$ has off-diagonal terms: galaxy 1 measures the combination $x_1 + x_2$ well
(eigenvalue 2500 along $(1,1)/\sqrt2$) and $x_1 - x_2$ poorly (eigenvalue 100). The components mix.
(d) Matrix weights are required; scalar weights mislabel the result. The same matrix weighting reappears in the
reference shear (B4) and in BA14 (G1).
}

% --------------------------------------------------------------------
\problem{A4}{Posterior, prior and likelihood}

\begin{given}
Bayes' rule: $p(x \mid d) = L(x)\,\pi(x) / Z$, with $L(x) = p(d\mid x)$ the likelihood, $\pi(x)$ the prior and $Z$
a normalising constant. A neural posterior estimator (NPE) trained on simulations whose parameters were drawn from
$\pitrain$ learns $q(x\mid d) \approx p(x\mid d)$ under that prior. If posterior samples were drawn under an
\emph{interim} prior, dividing by that prior (reweighting the samples) recovers the likelihood.
\source{\cHoggBayes; \cInterim.}
\end{given}

\need{nothing beyond the definitions.}

\begin{enumerate}
  \item Show that $q(x\mid d)/\pitrain(x) \propto L(x)$ wherever $\pitrain(x) > 0$.
  \item Our training prior is uniform on the box $|g_k| < 0.1$. What does $q/\pitrain$ look like inside the box, and
    what do we know about $L$ outside it?
  \item Why must we divide by the prior the network was \emph{trained} with, not by some other prior we prefer?
\end{enumerate}

\solution{
(a) $q = L\,\pitrain/Z$, so $q/\pitrain = L/Z \propto L$ (constant in $x$).
(b) Inside the box $\pitrain$ is constant, so $q/\pitrain \propto q$: the likelihood has the posterior's shape there.
Outside the box $q = 0$, so the likelihood is unknown there: the NPE says nothing about shears it never saw.
(c) The network has absorbed $\pitrain$ into $q$. Dividing by a different $\pi'$ leaves a factor $\pitrain/\pi'$
that is not part of the likelihood, which would bias every later product of likelihoods.
}

% --------------------------------------------------------------------
\problem{A5}{\py\ Why one galaxy's posterior mean is shrunk}

\begin{given}
Take one shear component. The likelihood is Gaussian, $L(x) \propto \exp[-(x-d)^2/2\sigma^2]$, with $\sigma = 0.07$
(the untruncated width of a typical KL likelihood), and the prior is uniform on $|x| < 0.1$. The posterior is then a
Gaussian truncated to $[-0.1, 0.1]$. Its mean is the ``posterior-mean shear estimate'' of one galaxy. The data scatter
as $d \sim \mathcal N(x_{\rm true}, \sigma^2)$. \emph{Derived here}.
\end{given}

\need{A4 (posterior = likelihood $\times$ prior).}

\begin{enumerate}
  \item Let $\alpha = (-0.1 - d)/\sigma$, $\beta = (0.1 - d)/\sigma$, and let $\phi$ and $\Phi$ be the standard normal
    density and CDF. Show that the posterior mean is
    $d + \sigma\,[\phi(\alpha) - \phi(\beta)]/[\Phi(\beta) - \Phi(\alpha)]$. (Integrate $x\,e^{-(x-d)^2/2\sigma^2}$ by
    substitution.)
  \item Explain without calculation why its average over $d$ is smaller in magnitude than $x_{\rm true}$.
  \item \py\ Compute $\E_d[\text{posterior mean}]$ for $x_{\rm true} \in [-0.05, 0.05]$, fit a line
    $(1+m)x_{\rm true} + c$, and report $m$.
  \item \py\ Now average $N$ independent galaxies at the same $x_{\rm true}$ by multiplying their likelihoods, which is
    a Gaussian of width $\sigma/\sqrt N$ by A2, truncated to the box. Repeat (c) for $N = 1, 2, 4, 8, 16$. How large
    must $N$ be for $|m| < 0.01$?
\end{enumerate}

\solution{
(a) With $u = (x-d)/\sigma$, $\int_{-0.1}^{0.1} x\,\phi(u)\,dx \propto d\,[\Phi(\beta)-\Phi(\alpha)] +
\sigma\int_\alpha^\beta u\,\phi(u)\,du$, and $\int_\alpha^\beta u\phi\,du = \phi(\alpha) - \phi(\beta)$. Dividing by the
normalisation $\Phi(\beta) - \Phi(\alpha)$ gives the result.
(b) When $d$ falls near or beyond an edge, the box cuts off the far tail of the likelihood, so the mean is dragged
inward. For $x_{\rm true} > 0$ the upper edge is reached more often than the lower one, so on average the estimate
is pulled toward zero.
(c) $m \approx -0.57$ for $\sigma = 0.07$. A single posterior mean is strongly shrunk, close to the $m \approx -0.6$
of the raw NPE posterior means (paper Sec.~6.1).
(d) $m \approx -0.57, -0.36, -0.17, -0.052, -0.0085$ for $N = 1, 2, 4, 8, 16$ (and $-0.0005$ for $N = 32$). The product
narrows as $\sigma/\sqrt N$ and stops touching the box, so the shrinkage disappears. With every galaxy at
$\sigma = 0.07$, $|m| < 0.01$ needs $N \approx 16$. Real KL widths vary from galaxy to galaxy, and the narrow ones
dominate a product, which is why the paper's stacks reach percent-level $m$ at smaller $N$. Multiplying likelihoods,
rather than averaging posterior means, is what removes the bias.
}

\section*{Part B. Estimators and stacking}\addcontentsline{toc}{section}{Part B. Estimators and stacking}
% ============================== Part B ==============================
\problem{B1}{Bias, scatter and pulls}

\begin{given}
An estimator $\hat\theta$ of a true value $\theta$ comes with a reported error bar $\hat\sigma$. The \emph{pull} is
$(\hat\theta - \theta)/\hat\sigma$. Suppose $\hat\theta \sim \mathcal N(\theta + b, s^2)$: bias $b$, true scatter $s$.
\emph{Derived here}; the pull is a standard diagnostic.
\end{given}

\need{nothing beyond Gaussians.}

\begin{enumerate}
  \item With $\hat\sigma$ fixed, show that the pull is $\mathcal N(b/\hat\sigma,\ s^2/\hat\sigma^2)$.
  \item What do ``pull mean $\neq 0$'' and ``pull spread $\neq 1$'' each tell you?
  \item Our 50-realization ensemble gives a pull spread of 0.84. By what factor do the reported error bars exceed
    the true scatter? If you used them in an inference, what fraction of the available information
    ($\propto 1/\sigma^2$) would you be claiming?
  \item The polar cells give 1437 pulls. What is the statistical uncertainty on their mean? Is a pull mean of
    $-0.02$ evidence of bias?
\end{enumerate}

\solution{
(a) $(\hat\theta - \theta)/\hat\sigma$ is a linear function of a Gaussian, with mean $b/\hat\sigma$ and standard
deviation $s/\hat\sigma$.
(b) A non-zero mean means bias, in units of the error bar. A spread above 1 means the error bars are too small
(overconfident); below 1, too large (conservative).
(c) $\hat\sigma/s = 1/0.84 = 1.19$, so the errors are 19\% too large. You would claim
$s^2/\hat\sigma^2 = 0.84^2 = 0.71$ of the information, as if you had thrown away 29\% of the galaxies. This loses
precision but introduces no bias.
(d) $0.84/\sqrt{1437} = 0.022$. A mean of $-0.02$ is $0.9\sigma$: consistent with zero.
}

% --------------------------------------------------------------------
\problem{B2}{Multiplicative and additive bias; shrinkage of a posterior mean}

\begin{given}
Shear biases are summarised by a linear response $\hat g = (1+m)\,g + c$, with $m$ the multiplicative and $c$ the
additive bias \source{\cBAresponse.} Model one galaxy with a Gaussian likelihood $d \sim \mathcal N(g, \sigma^2)$ and a
Gaussian prior $g \sim \mathcal N(0, \sigma_\pi^2)$. \emph{Derived here} (it reproduces the toy model of paper
Sec.~6.2).
\end{given}

\need{A1 (product of two Gaussians).}

\begin{enumerate}
  \item Treat the prior as a second ``measurement'' $0 \pm \sigma_\pi$. Show that the posterior mean is
    $\hat g = d\,\sigma_\pi^2/(\sigma_\pi^2 + \sigma^2)$.
  \item Average over the noise to find $\E[\hat g]$ and read off $m$ and $c$.
  \item Use $\sigma = 0.07$ and the standard deviation of a uniform prior on $[-0.1, 0.1]$,
    $\sigma_\pi = 0.2/\sqrt{12}$, to estimate $m$. Compare with A5(c).
\end{enumerate}

\solution{
(a) By A1, with weights $1/\sigma^2$ on $d$ and $1/\sigma_\pi^2$ on $0$:
$\hat g = (d/\sigma^2)/(1/\sigma^2 + 1/\sigma_\pi^2) = d\,\sigma_\pi^2/(\sigma_\pi^2+\sigma^2)$.
(b) $\E[d] = g$, so $\E[\hat g] = g\,\sigma_\pi^2/(\sigma_\pi^2 + \sigma^2)$: $m = -\sigma^2/(\sigma^2 + \sigma_\pi^2)$
and $c = 0$.
(c) $\sigma_\pi = 0.0577$, so $m = -0.0049/(0.0049 + 0.00333) = -0.60$. The truncated box of A5 gave $-0.57$; the
Gaussian toy captures the effect well. Shrinkage is strongest when the likelihood is wide compared with the prior.
}

% --------------------------------------------------------------------
\problem{B3}{Stack likelihoods, not posteriors}

\begin{given}
Each galaxy's NPE posterior is $q_i(g) \propto L_i(g)\,\pitrain(g)$ (A4). $N$ galaxies share one $g$ and have
independent data. The correct joint posterior is $p(g\mid d_1,\dots,d_N) \propto \pi(g)\prod_i L_i(g)$, and samples
drawn under an interim prior are combined by dividing that prior out.
\source{\cBAconst; \cInterim.}
\end{given}

\need{A4, B2.}

\begin{enumerate}
  \item Show that $\prod_i q_i(g) \propto \pitrain(g)^N\prod_i L_i(g)$. How many times is the prior counted?
  \item Use the Gaussian model of B2 (prior $\mathcal N(0,\sigma_\pi^2)$, each likelihood width $\sigma$). Show that the
    mean of $\prod_i q_i$ is shrunk by the same factor $\sigma_\pi^2/(\sigma_\pi^2+\sigma^2)$ for every $N$.
  \item What is the mean of the correct posterior $\pi\prod_i L_i$, and how does its $m$ behave as $N \to \infty$?
  \item Why does our stack use $L_i = q_i/\pitrain$ and a single flat prior on the cell shear?
\end{enumerate}

\solution{
(a) $\prod_i q_i \propto \prod_i L_i\pitrain = \pitrain^N\prod_i L_i$: the prior is counted $N$ times.
(b) $\pitrain^N$ has precision $N/\sigma_\pi^2$ and $\prod L_i$ has precision $N/\sigma^2$ about the mean $\bar d$, so
the mean is $\bar d\,(N/\sigma^2)/(N/\sigma^2 + N/\sigma_\pi^2) = \bar d\,\sigma_\pi^2/(\sigma_\pi^2+\sigma^2)$: the
$N$ cancels, and $m = -0.6$ no matter how many galaxies you multiply.
(c) $\bar d\,(N/\sigma^2)/(N/\sigma^2 + 1/\sigma_\pi^2)$, so $m = -\sigma^2/(\sigma^2 + N\sigma_\pi^2) \to 0$ as
$N \to \infty$.
(d) Dividing out $\pitrain$ leaves the likelihoods, which multiply correctly. The cell then has one prior (uniform
on $|\vG| < 0.1$), counted once, and the data dominate as $N$ grows.
}

% --------------------------------------------------------------------
\problem{B4}{Fitting one shear to galaxies whose shears differ: the reference shear}

\begin{given}
Fitting a model that cannot be exactly right (a misspecified model) converges, as the data grow, to the parameter
value that makes the model's distribution closest to the true one in Kullback--Leibler divergence: the
\emph{pseudo-true} value.
\source{\cMisspec.}
In a cell, galaxy $i$ has true shear $\vt_i$ (in the cell frame), a Gaussian likelihood centred on $\vd_i$ with
precision $F_i$, and $\E[\vd_i] = \vt_i$. The stack fits one common $\vG$ to all of them.
\end{given}

\need{A1, A3 (the matrix-weighted mean).}

\begin{enumerate}
  \item By A3, $\hat\vG = \Ftot^{-1}\sum_i F_i\vd_i$. Average over the noise, assuming $F_i$ doesn't depend on it.
    Show that $\E[\hat\vG] = \Ftot^{-1}\sum_i F_i\vt_i \equiv \gref$.
  \item Show that $\mathrm{Cov}(\hat\vG) = \Ftot^{-1}$, and that it doesn't depend on how much the $\vt_i$ differ.
    Why is this still the right error bar \emph{about} $\gref$?
  \item Scalar example: galaxy 1 has $t_1 = 0.02$, $\sigma_1 = 0.02$; galaxy 2 has $t_2 = 0$, $\sigma_2 = 0.06$.
    Compute $g_{\rm ref}$ and the plain mean. Which one does the stack estimate?
  \item Why would comparing the stack with the shear at the cell centre be the wrong test?
\end{enumerate}

\solution{
(a) $\E[\hat\vG] = \Ftot^{-1}\sum_i F_i\,\E[\vd_i] = \Ftot^{-1}\sum_i F_i\vt_i$. This is the pseudo-true value for
Gaussian likelihoods: the constant that best describes the varying $\vt_i$, with each galaxy weighted by its
precision.
(b) With independent noises of covariance $F_i^{-1}$:
$\mathrm{Cov}(\hat\vG) = \Ftot^{-1}\big(\sum_i F_iF_i^{-1}F_i\big)\Ftot^{-1} = \Ftot^{-1}$. The $\vt_i$ don't appear.
It is the right error about $\gref$ because $\gref$ is exactly what $\hat\vG$ estimates; so
$\hat\vG - \gref \sim \mathcal N(0, \Ftot^{-1})$, and the pulls about $\gref$ are standard normal.
(c) $w_1 = 2500$, $w_2 = 278$: $g_{\rm ref} = 50/2778 = 0.0180$; plain mean $0.01$. The stack estimates $0.0180$,
with error $0.019$. Comparing it with $0.01$ would call a correct stack biased.
(d) The cell-centre shear is not what the stack converges to. The difference is a deterministic binning effect,
largest where the shear changes quickly. (Full derivation: \texttt{CELL\_REFERENCE\_SHEAR.md}, Secs.~3--5.)
}

% --------------------------------------------------------------------
\problem{B5}{\py\ The tangential frame, and checking a stack with pulls}

\begin{given}
About a lens at the origin, a galaxy at position angle $\phi$ has tangential and cross shear
$g_t = -\mathrm{Re}(g\,e^{-2i\phi})$, $g_\times = -\mathrm{Im}(g\,e^{-2i\phi})$, with $g = g_1 + i g_2$.
\source{\cBStan.}
As a matrix, $(g_t, g_\times)^\T = R(\phi)\,(g_1, g_2)^\T$ with
$R(\phi) = -\begin{pmatrix}\cos2\phi & \sin2\phi\\ -\sin2\phi & \cos2\phi\end{pmatrix}$.
\end{given}

\need{A3, B1, B4.}

\begin{enumerate}
  \item Show that $R(\phi)$ is orthogonal ($RR^\T = I$), and that a likelihood with precision $F$ in the image frame has
    precision $F' = RFR^\T$ in the tangential frame.
  \item An axisymmetric lens produces purely tangential shear: $g_\times = 0$ and $g_t > 0$. Write $(g_1, g_2)$ at angle
    $\phi$ in terms of $g_t$.
  \item \py\ Simulate a cell of 5 galaxies at different $\phi$ around a lens with $g_t$ varying from 0.01 to 0.03 and
    anisotropic, randomly oriented likelihood covariances. Draw noisy measurements, rotate each into the tangential
    frame, stack with A3, and repeat 2000 times. Show that the pulls about $\gref$ have mean $\approx 0$ and spread
    $\approx 1$, while the pulls about the plain mean of the $g_{t,i}$ are biased.
\end{enumerate}

\solution{
(a) Writing $R = -Q$ with $Q$ a rotation by $2\phi$, $RR^\T = QQ^\T = I$. If
$\ln L \propto -\tfrac12 (\vg - \hat\vg)^\T F(\vg - \hat\vg)$ and $\vg = R^\T\vg'$, then
$(\vg - \hat\vg)^\T F(\vg - \hat\vg) = (\vg' - R\hat\vg)^\T RFR^\T(\vg' - R\hat\vg)$, so $F' = RFR^\T$.
(b) $g\,e^{-2i\phi} = -g_t$ gives $g = -g_t e^{2i\phi}$: $g_1 = -g_t\cos2\phi$, $g_2 = -g_t\sin2\phi$.
(c) See the solutions notebook: the pulls about $\gref$ have mean and spread close to $0$ and $1$. The pulls about
the plain mean are offset (by about $0.05$--$0.1\sigma$ in this example; more when the precisions differ more), because
precise galaxies at particular $g_t$ dominate the stack.
}

\section*{Part C. Weak-lensing basics}\addcontentsline{toc}{section}{Part C. Weak-lensing basics}
% ============================== Part C ==============================
\problem{C1}{One potential gives both mass and shear}

\begin{given}
A thin lens is described by the lensing potential $\psi(\bm\theta)$. Its second derivatives
($\psi_{ij} = \partial^2\psi/\partial\theta_i\partial\theta_j$) give the convergence (projected mass in units of the
critical density) and the two shear components:
\[
\kappa = \tfrac12(\psi_{11} + \psi_{22}),\qquad
\gamma_1 = \tfrac12(\psi_{11} - \psi_{22}),\qquad
\gamma_2 = \psi_{12}.
\]
\source{\cBSpsi.}
A singular isothermal sphere (SIS) has $\psi = \theta_E\,|\bm\theta|$, with Einstein radius $\theta_E$.
\source{\cBSsis.}
\end{given}

\need{nothing (calculus only).}

\begin{enumerate}
  \item Compute $\psi_{11}$, $\psi_{22}$ and $\psi_{12}$ for the SIS. Write $\theta_1 = \theta\cos\phi$,
    $\theta_2 = \theta\sin\phi$.
  \item Show that $\kappa = \theta_E/(2\theta)$ and that $\gamma \equiv \gamma_1 + i\gamma_2 = -\frac{\theta_E}{2\theta}\,e^{2i\phi}$.
  \item Why does knowing $\kappa$ everywhere fix $\gamma$ everywhere (up to what)?
\end{enumerate}

\solution{
(a) $\partial_1\psi = \theta_E\theta_1/\theta$. Then $\psi_{11} = \theta_E\theta_2^2/\theta^3$,
$\psi_{22} = \theta_E\theta_1^2/\theta^3$ and $\psi_{12} = -\theta_E\theta_1\theta_2/\theta^3$.
(b) $\kappa = \tfrac12\theta_E(\theta_1^2 + \theta_2^2)/\theta^3 = \theta_E/(2\theta)$.
$\gamma_1 = \theta_E(\theta_2^2 - \theta_1^2)/(2\theta^3) = -\tfrac{\theta_E}{2\theta}\cos2\phi$ and
$\gamma_2 = -\tfrac{\theta_E}{2\theta}\sin2\phi$, so $\gamma = -\tfrac{\theta_E}{2\theta}e^{2i\phi}$.
(c) $\kappa$ fixes $\nabla^2\psi$, which fixes $\psi$ up to a harmonic function (one with $\nabla^2 = 0$, such as a
constant gradient or a term $\propto \theta^2$ that is really a uniform sheet; see C3). The shear follows from $\psi$.
This link is the basis of every mass map (Part D).
}

% --------------------------------------------------------------------
\problem{C2}{Tangential and cross shear}

\begin{given}
About the lens centre, the tangential and cross components of the shear are
$\gamma_t = -\mathrm{Re}(\gamma\,e^{-2i\phi})$ and $\gamma_\times = -\mathrm{Im}(\gamma\,e^{-2i\phi})$.
\source{\cBStan.}
For any mass distribution, the tangential shear averaged around a circle of radius $\theta$ is
$\langle\gamma_t\rangle(\theta) = \bar\kappa(\le\theta) - \langle\kappa\rangle(\theta)$, where $\bar\kappa(\le\theta)$ is
the mean convergence inside the circle and $\langle\kappa\rangle(\theta)$ the mean on it. For an axially symmetric lens
the averages can be dropped.
\source{\cBSgtmean.}
\end{given}

\need{C1 (SIS $\kappa$ and $\gamma$).}

\begin{enumerate}
  \item Use C1(b) to show that the SIS has $\gamma_t = \theta_E/(2\theta)$ and $\gamma_\times = 0$.
  \item Check the result against $\bar\kappa(<\theta) - \kappa(\theta)$.
  \item A subhalo next to a main halo breaks the axial symmetry about the main-halo centre. What does it add to
    $\gamma_\times$ measured about the main centre, and why do our polar cells stack in $(\gamma_t, \gamma_\times)$?
\end{enumerate}

\solution{
(a) $\gamma e^{-2i\phi} = -\tfrac{\theta_E}{2\theta}$ is real and negative, so $\gamma_t = \theta_E/(2\theta)$ and
$\gamma_\times = 0$.
(b) $\bar\kappa(<\theta) = \frac{1}{\pi\theta^2}\int_0^\theta \frac{\theta_E}{2\theta'}\,2\pi\theta'\,d\theta' = \theta_E/\theta$,
so $\bar\kappa - \kappa = \theta_E/\theta - \theta_E/(2\theta) = \theta_E/(2\theta)$. It agrees.
(c) The subhalo's shear is tangential about \emph{its own} centre, so about the main centre it has both $\gamma_t$ and
$\gamma_\times$ parts. A non-zero $\gamma_\times$ about the main centre signals non-axisymmetric structure. In the
$(\gamma_t, \gamma_\times)$ frame the main halo's signal sits in one component, so departures from it stand out.
}

% --------------------------------------------------------------------
\problem{C3}{Reduced shear and the mass-sheet degeneracy}

\begin{given}
Galaxy shapes respond to the reduced shear $g = \gamma/(1-\kappa)$, not to $\gamma$.
\source{\cBSreduced.}
The transformation $\kappa \to \lambda\kappa + (1-\lambda)$, $\gamma \to \lambda\gamma$ (adding a uniform sheet of
mass and rescaling) leaves $g$ unchanged: the mass-sheet degeneracy.
\source{\cBSmsd.}
\end{given}

\need{C1.}

\begin{enumerate}
  \item Verify that $g$ is invariant under the transformation.
  \item For $\kappa = 0.05$ and $\gamma = 0.02$, compute $g$. How large is the difference from $\gamma$, in percent?
  \item What can shear data \emph{not} tell you about a $\kappa$ map?
\end{enumerate}

\solution{
(a) $g' = \lambda\gamma/(1 - \lambda\kappa - 1 + \lambda) = \lambda\gamma/[\lambda(1-\kappa)] = g$.
(b) $g = 0.02/0.95 = 0.02105$, 5\% above $\gamma$. That is negligible far out, but not near the cluster centre, where
$\kappa \sim 0.2$.
(c) The overall level (and scale) of $\kappa$: only its shape is fixed. A mass map's zero point is a convention.
}

% --------------------------------------------------------------------
\problem{C4}{Adding two halos}

\begin{given}
$\kappa$ is the projected mass, and $\psi$ obeys the linear Poisson equation $\nabla^2\psi = 2\kappa$. So the $\psi$,
$\kappa$ and $\gamma$ of two halos are the sums of their separate values.
\source{\cBSpoisson.}
\end{given}

\need{C1, C3.}

\begin{enumerate}
  \item Halo A gives $\kappa_A = 0.05$, $\gamma_A = (0.02, 0)$ at some point; halo B gives $\kappa_B = 0.03$,
    $\gamma_B = (0, 0.01)$. Compute the correct reduced shear of A+B.
  \item Compute the incorrect $g_A + g_B$. How big is the error?
  \item Why do our simulations add $\gamma$ and $\kappa$ over the halos first and form $g$ only at the end?
\end{enumerate}

\solution{
(a) $\gamma = (0.02, 0.01)$ and $\kappa = 0.08$: $g = (0.02174, 0.01087)$.
(b) \[g_A + g_B = (0.02/0.95,\ 0.01/0.97) = (0.02105,\ 0.01031).\] The error is $\approx 3$--5\%, as large as the effects
we test for.
(c) Because $\gamma$ and $\kappa$ superpose and $g$ does not.
}

% --------------------------------------------------------------------
\problem{C5}{Where the subhalo's signal is, and how strong it is}

\begin{given}
In our simulation (subhalo of $3\times10^{14}\,M_\odot$ at $z = 0.13$), the subhalo's extra reduced shear at distance $r$
from it is close to $|\Delta g| \approx A/r$ with $A \approx 0.009$\,Mpc. That is SIS-like; the true profile is NFW,
whose shear formulae are in \cWBnfw. The sources have density $n \approx 5\,{\rm Mpc}^{-2}$, and each measures each
shear component with noise $\sigma \approx 0.05$. For a signal of known shape, the best possible (matched-filter)
significance is $({\rm S/N})^2 = \sum_i |\Delta\vg_i|^2/\sigma^2$. \emph{Derived here} (numbers from our
\texttt{s3\_sub1} manifest).
\end{given}

\need{A2 (information adds).}

\begin{enumerate}
  \item Replace the sum by an integral over a uniform source density and show that
    $({\rm S/N})^2 = 2\pi n A^2\sigma^{-2}\ln(r_{\max}/r_{\min})$.
  \item Evaluate it for $r_{\min} = 0.1$\,Mpc and $r_{\max} = 3$\,Mpc.
  \item What fraction of $({\rm S/N})^2$ comes from within 1\,Mpc, and from within 0.5\,Mpc, of the subhalo?
  \item With each galaxy's full, anisotropic precision matrix in place of an isotropic $\sigma$, the same calculation
    gives ${\rm S/N} \approx 3.2$ for this galaxy set. Why is it larger?
\end{enumerate}

\solution{
(a) $\sum_i \to \int n\,2\pi r\,dr$: $({\rm S/N})^2 = \int_{r_{\min}}^{r_{\max}} n\,2\pi r\,\frac{A^2}{r^2\sigma^2}\,dr
= 2\pi nA^2\sigma^{-2}\ln(r_{\max}/r_{\min})$.
(b) $A^2/\sigma^2 = 0.0324$, so $2\pi\cdot5\cdot0.0324 = 1.02$, times $\ln 30 = 3.40$: $({\rm S/N})^2 = 3.46$, so
${\rm S/N} \approx 1.9$. (The actual galaxies give 1.65 for isotropic $\sigma = 0.049$.)
(c) Each factor of $e$ in radius contributes equally: within 1\,Mpc, $\ln 10/\ln 30 = 0.68$; within 0.5\,Mpc,
$\ln 5/\ln 30 = 0.47$. (The simulation gives 0.58 and 0.35; the NFW profile falls faster than $1/r$ at large $r$
and is shallower near the centre.) The signal is spread over a few Mpc, not concentrated in one spot, which is why a
method that uses every galaxy beats one cell.
(d) KL measures one shear direction much better than the other (median $\sigma \approx 0.039$ along the best direction
vs $0.057$ along the worst). With matrix weights, each galaxy contributes most where it is precise (A3). This is the
benchmark for Part G: no map can beat the matched filter.
}

\section*{Part D. Fourier space and Kaiser--Squires}\addcontentsline{toc}{section}{Part D. Fourier space and Kaiser--Squires}
% ============================== Part D ==============================
\problem{D1}{Derivatives become multiplications}

\begin{given}
With the Fourier transform $\hat f(\vk) = \int f(\bm\theta)\,e^{i\vk\cdot\bm\theta}\,d^2\theta$ (the sign convention of
B\&S01, Sec.~5.2.1, Eq.~(5.1)), a derivative $\partial/\partial\theta_j$ of $f$ becomes a factor $\mp i k_j$ on
$\hat f$. The sign depends on the convention and will cancel in everything below, because every quantity uses
\emph{two} derivatives. The potential obeys the Poisson equation $\nabla^2\psi = 2\kappa$.
\source{\cBSpoisson.}
\end{given}

\need{C1 ($\kappa$ and $\gamma$ as second derivatives of $\psi$).}

\begin{enumerate}
  \item Show that $\widehat{\partial_i\partial_j f} = -k_ik_j\hat f$ for either sign convention.
  \item Show that $\hat\psi = -2\hat\kappa/k^2$ for $\vk \neq 0$.
  \item What happens at $\vk = 0$? Relate this to C3.
\end{enumerate}

\solution{
(a) Each derivative gives $\mp ik$; two give $(\mp i)^2k_ik_j = -k_ik_j$.
(b) $\nabla^2 \to -(k_1^2 + k_2^2) = -k^2$, so $-k^2\hat\psi = 2\hat\kappa$.
(c) The equation says nothing about $\hat\kappa(0)$, the mean $\kappa$ over the field: shear never constrains it. It is
the mass-sheet degeneracy of C3 in Fourier form.
}

% --------------------------------------------------------------------
\problem{D2}{The Kaiser--Squires kernel}

\begin{given}
The shear is a convolution of $\kappa$ with a kernel. In Fourier space, one multiplication per mode links them, and
the inversion recovers $\hat\kappa$ from $\hat\gamma$ by the complex-conjugate kernel. The papers write the kernel as
$\hat D = \pi(k_1^2 - k_2^2 + 2ik_1k_2)/k^2$ with a factor $\pi^{-1}$ in front of the product; below we absorb the
$\pi$.
\source{\cBSks; \cJeffreyKS.}
\end{given}

\need{C1, D1.}

\begin{enumerate}
  \item Using C1 and D1, show that $\hat\gamma = D\,\hat\kappa$ with $\hat\gamma = \hat\gamma_1 + i\hat\gamma_2$ and
    $D(\vk) = (k_1^2 - k_2^2 + 2ik_1k_2)/k^2$.
  \item Show that $|D| = 1$. Write $D = e^{2i\beta}$, with $\beta$ the direction of $\vk$.
  \item Show that $\hat\kappa = D^*\hat\gamma$ for $\vk \neq 0$. Why is this inversion exact for noise-free data on an
    infinite field, yet noisy in practice?
\end{enumerate}

\solution{
(a) $\hat\gamma_1 = \tfrac12(\widehat{\psi_{11}} - \widehat{\psi_{22}}) = \tfrac12(-k_1^2 + k_2^2)\hat\psi
= \frac{k_1^2 - k_2^2}{k^2}\hat\kappa$ and $\hat\gamma_2 = \widehat{\psi_{12}} = -k_1k_2\hat\psi = \frac{2k_1k_2}{k^2}\hat\kappa$.
(b) $|D|^2 = [(k_1^2 - k_2^2)^2 + 4k_1^2k_2^2]/k^4 = (k_1^2 + k_2^2)^2/k^4 = 1$. With $\vk = k(\cos\beta, \sin\beta)$,
$D = \cos2\beta + i\sin2\beta = e^{2i\beta}$.
(c) $D^*D = |D|^2 = 1$. Noise $\hat n$ in the shear map passes into $\hat\kappa$ undamped ($|D^*| = 1$) at every
scale, including small scales where the signal is tiny. KS therefore needs smoothing afterwards (D4).
}

% --------------------------------------------------------------------
\problem{D3}{Finite fields, padding and the mean mass}

\begin{given}
The inversion holds for $\vk \neq 0$; the real-space form carries an undetermined additive constant $\kappa_0$, because
a uniform surface mass density produces no shear.
\source{\cBSks; \cBSmsd.}
On a grid, the discrete Fourier transform treats the field as periodic.
\end{given}

\need{D1, D2.}

\begin{enumerate}
  \item Why does a periodic transform make the left edge of a shear map act on the right edge? What does padding the
    map with zeros to twice its size do?
  \item Our KS map sets $\hat\kappa(0) = 0$. What does that assume, and can the data test it?
  \item Our likelihood $\kappa$ map pads with zeros too, in the forward direction ($\kappa \to \gamma$). What
    assumption does that make about $\kappa$ outside the field?
\end{enumerate}

\solution{
(a) In a periodic box, mass near one edge is also ``next to'' the opposite edge, and the non-local kernel (shear
falls only as $1/r^2$) couples them. Padding with zeros places the copies far away and removes most of that
wrap-around.
(b) That the mean $\kappa$ in the field is zero. The data cannot test it (D1(c)), so the map's zero point is a
convention.
(c) That $\kappa = 0$ outside the field. Mass outside still shears galaxies inside, so this is an approximation that
matters most near the edge.
}

% --------------------------------------------------------------------
\problem{D4}{\py\ What binning does to a tangential pattern}

\begin{given}
A lens at the origin produces purely tangential shear, $g = -g_t(r)\,e^{2i\phi}$ (B5(b)). KS needs the shear on a
regular grid, so it averages the galaxies in each cell. Smoothing by a Gaussian of width $s$ multiplies each Fourier
mode by $e^{-k^2s^2/2}$ (the convolution theorem). \emph{Derived here}; the cell-averaging argument is ours. Our cells
are 0.98\,Mpc squares and our subhalo sits almost at a cell centre.
\end{given}

\need{B5(b), D2.}

\begin{enumerate}
  \item Show that the average of $g_1$ and $g_2$ over a square centred on the lens is zero for any $g_t(r)$.
    (Use the symmetry of $\cos2\phi$ and $\sin2\phi$ under $\phi \to \phi + \pi/2$ and $\phi \to -\phi$.)
  \item So what does the KS cell containing the subhalo report about the subhalo? Where does KS's information about
    the subhalo come from?
  \item \py\ Compute the cell-averaged shear of an SIS-like subhalo ($g_t = A/r$) for cells whose centre is offset from
    the subhalo by 0 to 1 cell, and compare with the subhalo's shear at the cell centre.
  \item Gaussian smoothing with $s = 0.7$ cell multiplies the mode at the grid's Nyquist scale ($k = \pi/\Delta$, cell
    size $\Delta$) by what factor? What does that mean for the map's resolution?
\end{enumerate}

\solution{
(a) The square is symmetric under a $90^\circ$ rotation, which sends $\cos2\phi \to -\cos2\phi$ and
$\sin2\phi \to -\sin2\phi$ while $r$ is unchanged, so each average equals minus itself.
(b) Nothing (to first order): the cell centred on the subhalo averages its signal away. KS learns about the subhalo
only from the neighbouring cells, where the pattern no longer cancels. The likelihood map, which compares each galaxy
with the shear at its own position, keeps that signal.
(c) See the notebook: the cell average is zero at offset 0 and grows to about the centre value once the cell is
about one cell width away. The nearest cells are where the cancellation hurts.
(d) $e^{-(\pi\cdot0.7)^2/2} = e^{-2.42} = 0.089$. Structure at the scale of one or two cells is suppressed about
tenfold, so KS resolves only features several cells (several Mpc) across.
}

\section*{Part E. Regularised inversion and the Wiener filter}\addcontentsline{toc}{section}{Part E. Regularised inversion and the Wiener filter}
% ============================== Part E ==============================
\problem{E1}{When the data can't determine the answer}

\begin{given}
A linear model $\vd = A\mathbf{s} + \mathbf n$, with unknowns $\mathbf s$ and noise covariance $N$. Weighted least squares
gives $\hat{\mathbf s} = [A^\T N^{-1}A]^{-1}A^\T N^{-1}\vd$, which needs $A^\T N^{-1}A$ to be invertible.
\source{\cHBLwls.}
\end{given}

\need{A3 (the matrix least-squares solution).}

\begin{enumerate}
  \item One datum measures the sum of two unknowns: $d = s_1 + s_2 + n$, so $A = (1\ \ 1)$. Show that $A^\T A$ is
    singular. Describe all the solutions that fit the datum perfectly.
  \item Two data, $A = \mathrm{diag}(1, 0.01)$, $N = \sigma^2 I$. Show that $\hat s_2$ amplifies the noise 100-fold.
  \item Our $\kappa$ map has 144 unknown pixels and about 150 galaxies, each with noise $\approx 0.05$ against a subhalo
    signal of $\approx 0.01$. Which of (a) and (b) describes it?
\end{enumerate}

\solution{
(a) $A^\T A = \begin{pmatrix}1&1\\1&1\end{pmatrix}$ has determinant 0. Any $(s_1, s_2)$ with $s_1 + s_2 = d$ fits; the
data fix one combination and say nothing about $s_1 - s_2$.
(b) $\hat s_2 = d_2/0.01 = s_2 + 100\,n_2$: a direction the data barely respond to is recovered with 100 times the noise.
(c) Both. Small-scale patterns of $\kappa$ produce shears that nearly cancel at the galaxies (small singular values of
$A$, as in (b)), and the noisy, nearly degenerate system has many near-equally good solutions (as in (a)). Something
other than the data must choose between them: a prior.
}

% --------------------------------------------------------------------
\problem{E2}{A Gaussian prior makes the problem well posed}

\begin{given}
Add a Gaussian prior $\mathbf s \sim \mathcal N(0, S)$ to the model of E1. For a Gaussian signal and Gaussian noise, the
linear minimum-variance (Wiener) filter is $W = SA^\T(ASA^\T + N)^{-1}$, and it is also the maximum and the mean of the
Gaussian posterior.
\source{\cZaroubiWF; \cJeffreyWF.}
\end{given}

\need{A3 (completing the square in matrix form), E1.}

\begin{enumerate}
  \item Write $-2\ln p(\mathbf s\mid\vd) = (\vd - A\mathbf s)^\T N^{-1}(\vd - A\mathbf s) + \mathbf s^\T S^{-1}\mathbf s + \text{const}$.
    Complete the square to show that the posterior is Gaussian with precision
    $\mathsf H = S^{-1} + A^\T N^{-1}A$ and mean $\hat{\mathbf s} = \mathsf H^{-1}A^\T N^{-1}\vd$.
  \item Show that $\hat{\mathbf s} = W\vd$, using the identity
    $(S^{-1} + A^\T N^{-1}A)^{-1}A^\T N^{-1} = SA^\T(ASA^\T + N)^{-1}$. Verify the identity by multiplying both sides on
    the left by $S^{-1} + A^\T N^{-1}A$ and on the right by $ASA^\T + N$.
  \item Redo E1(a) with $S = \sigma_s^2 I$ and $N = \sigma_n^2$. Is $\mathsf H$ invertible now? What is $\hat{\mathbf s}$?
\end{enumerate}

\solution{
(a) The $\mathbf s$-dependent terms are $\mathbf s^\T(A^\T N^{-1}A + S^{-1})\mathbf s - 2\mathbf s^\T A^\T N^{-1}\vd$,
which equals $(\mathbf s - \hat{\mathbf s})^\T\mathsf H(\mathbf s - \hat{\mathbf s})$ + const with
$\hat{\mathbf s} = \mathsf H^{-1}A^\T N^{-1}\vd$. A Gaussian's mean and maximum coincide.
(b) Left side times the two factors: $A^\T N^{-1}(ASA^\T + N) = A^\T N^{-1}ASA^\T + A^\T$. Right side:
$(S^{-1} + A^\T N^{-1}A)SA^\T = A^\T + A^\T N^{-1}ASA^\T$. They agree.
(c) $\mathsf H = \sigma_s^{-2}I + \sigma_n^{-2}\begin{pmatrix}1&1\\1&1\end{pmatrix}$ has eigenvalues
$\sigma_s^{-2} + 2\sigma_n^{-2}$ and $\sigma_s^{-2}$, both positive, so it is invertible. By symmetry
$\hat s_1 = \hat s_2 = \frac{d\,\sigma_s^2}{2\sigma_s^2 + \sigma_n^2}$: the datum is split equally, shrunk toward zero,
and the undetermined direction $s_1 - s_2$ is set to its prior value, 0.
}

% --------------------------------------------------------------------
\problem{E3}{One mode at a time: what the Wiener filter keeps}

\begin{given}
If $S$ and $N$ are diagonal in the same basis, for example Fourier modes of a stationary field, the filter acts on
each mode separately and suppresses the noise by the ratio prior/(prior + noise). It is the linear estimate with the
smallest residual variance.
\source{\cZaroubiMSE.}
\end{given}

\need{E2.}

\begin{enumerate}
  \item One mode, $d = s + n$, with $s \sim \mathcal N(0, S)$ and $n \sim \mathcal N(0, N)$. For an estimate
    $\hat s = w\,d$, show that the mean squared error $\E[(\hat s - s)^2] = (1-w)^2S + w^2N$, and that it is minimised
    at $w = S/(S+N)$.
  \item For a \emph{fixed} true $s$, show that $\E[\hat s\mid s] = ws$. In what sense is the Wiener estimate biased?
  \item Large scales have $S \gg N$ and small scales $S \ll N$. Describe the filtered map.
  \item The subhalo in our likelihood map recovers about 40\% of its true peak value. What value of $w$ does that
    suggest for the scales it occupies?
\end{enumerate}

\solution{
(a) $\hat s - s = (w-1)s + wn$; with $s$ and $n$ independent, the variance is $(1-w)^2S + w^2N$. Setting the
derivative to zero, $-2(1-w)S + 2wN = 0$, gives $w = S/(S+N)$.
(b) $\E[w(s+n)\mid s] = ws$: every mode is shrunk toward zero by $w < 1$. The estimate is biased for any given field,
but over the ensemble of fields its error is minimal. This is the trade made on purpose.
(c) Large scales pass almost unchanged and small scales are suppressed, so the map is a smoothed, slightly shrunk
version of the truth.
(d) $w \approx 0.4$: at the subhalo's scale the signal and noise powers are comparable ($S/N \approx 0.7$).
}

% --------------------------------------------------------------------
\problem{E4}{A smoothness penalty is a power spectrum}

\begin{given}
Our map prior is
$-\ln p(\kappa) = \frac{1}{2\sk^2}\big[\sum_j\kappa_j^2 + \ell^2\sum_{\langle jj'\rangle}(\kappa_j - \kappa_{j'})^2\big]$,
summed over pixels $j$ and neighbouring pairs $\langle jj'\rangle$. \emph{Derived here}; it is the regulariser of
\texttt{likelihood\_kappa} in \texttt{h\_cluster\_analysis.py}.
\end{given}

\need{E2, E3 (a Gaussian prior with covariance $S$; per-mode filtering).}

\begin{enumerate}
  \item On a periodic 1D chain of $M$ pixels, insert a Fourier mode $\kappa_j = a\,e^{ikj}$. Show that
    $\sum_j |\kappa_{j+1} - \kappa_j|^2 = 4\sin^2(k/2)\sum_j|\kappa_j|^2$.
  \item Conclude that each mode has prior variance $P(k) = \sk^2/[1 + 4\ell^2\sin^2(k/2)] \approx \sk^2/(1+\ell^2k^2)$
    for small $k$.
  \item What do $\sk$ and $\ell$ each control? Which scales does the prior suppress?
\end{enumerate}

\solution{
(a) $\kappa_{j+1} - \kappa_j = a\,e^{ikj}(e^{ik} - 1)$ and $|e^{ik} - 1|^2 = 2 - 2\cos k = 4\sin^2(k/2)$.
(b) Modes are independent (they diagonalise the quadratic form), and mode $k$ has precision
$[1 + 4\ell^2\sin^2(k/2)]/\sk^2$, whose inverse is the prior variance. For small $k$, $4\sin^2(k/2) \approx k^2$.
(c) $\sk$ sets the overall amplitude the prior allows. $\ell$ sets the scale below which fluctuations are suppressed:
the variance falls by half at $k \approx 1/\ell$. Combined with E3, the Wiener weight per mode is
$P(k)/[P(k) + N(k)]$, which removes small scales first.
}

% --------------------------------------------------------------------
\problem{E5}{\py\ A one-dimensional Wiener filter}

\begin{given}
A toy 1D ``$\kappa$'' made of a broad bump (the main halo) plus a narrow bump (the subhalo) is sampled by noisy point
measurements at random positions, each with noise $\sigma = 0.05$ of a quantity linear in $\kappa$. The MAP map under
the prior of E4 is $\hat\kappa = \mathsf H^{-1}A^\T N^{-1}\vd$ (E2). \emph{Derived here}.
\end{given}

\need{E2, E4.}

\begin{enumerate}
  \item \py\ Build $A$ (a measurement samples the pixel it falls in) and the prior precision of E4. Compute the MAP
    map for $\sk \in \{0.03, 0.1, 0.3\}$ and $\ell \in \{0.5, 1, 2\}$ pixels.
  \item \py\ Which settings recover the narrow bump best? What happens to it for small $\sk$, and for large $\ell$?
  \item \py\ Repeat with 20 noise realizations. For one setting, plot the mean map against the truth. Is the mean map
    equal to the truth?
\end{enumerate}

\solution{
(a)--(b) See the solutions notebook. Small $\sk$ shrinks everything (low peaks); large $\ell$ blurs the narrow bump into
the broad one. Large $\sk$ with small $\ell$ keeps the bump's height but lets the noise through. The best setting
balances the two, just as the tuning did on the real map.
(c) No: the mean map is a smoothed, shrunk version of the truth (E3(b)). Averaging over noise removes the noise, not
the prior's bias. Part F makes this precise with the resolution matrix.
}

\section*{Part F. Uncertainty, resolution and effective parameters}\addcontentsline{toc}{section}{Part F. Uncertainty, resolution and effective parameters}
% ============================== Part F ==============================
\problem{F1}{The Laplace approximation}

\begin{given}
Near its maximum $\mathbf x_0$, a log density is approximately quadratic,
$\ln P(\mathbf x) \simeq \ln P(\mathbf x_0) - \tfrac12(\mathbf x - \mathbf x_0)^\T\mathsf A(\mathbf x - \mathbf x_0)$, with
$\mathsf A = -\nabla\nabla\ln P|_{\mathbf x_0}$. Treating $P$ as a Gaussian with covariance $\mathsf A^{-1}$ is the
Laplace approximation.
\source{\cMacKayLaplace.}
\end{given}

\need{A1 (a Gaussian's log is a quadratic).}

\begin{enumerate}
  \item For $P(x) \propto x^a e^{-bx}$ ($x > 0$), find the mode, $\mathsf A$, and the Laplace standard deviation.
  \item The exact standard deviation is $\sqrt{a+1}/b$. Compare for $a = 4$. When does Laplace work well?
  \item For the posterior of E2, show that the Laplace approximation is exact, with $\mathsf A = \mathsf H$.
\end{enumerate}

\solution{
(a) $\ln P = a\ln x - bx$; the mode is $x_0 = a/b$, $\mathsf A = a/x_0^2 = b^2/a$, so $\sigma_L = \sqrt a/b$.
(b) $a = 4$: $2/b$ against $\sqrt5/b = 2.24/b$, 11\% low. It works well when the distribution is close to Gaussian,
here for large $a$ (the ratio $\sqrt{a/(a+1)} \to 1$).
(c) The log posterior is exactly quadratic, with $-\nabla\nabla\ln p = S^{-1} + A^\T N^{-1}A = \mathsf H$.
}

% --------------------------------------------------------------------
\problem{F2}{Gauss--Newton: curvature from the model, not from the data grid}

\begin{given}
For a least-squares objective $F(\mathbf x) = \tfrac12\sum_i f_i(\mathbf x)^2$, the exact Hessian is
$J^\T J + \sum_i f_i\,\nabla\nabla f_i$, with $J$ the Jacobian of the residuals $f_i$. Gauss--Newton keeps only $J^\T J$.
\source{\cGaussNewton.}
Our galaxies enter through $\ln L_i(\vg_i(\kappa))$, read off a grid by bilinear interpolation, and each $L_i$ has
precision $F_i$ (A3).
\end{given}

\need{A3, F1.}

\begin{enumerate}
  \item Suppose $\ln L_i(\vg) \approx -\tfrac12(\vg - \hat\vg_i)^\T F_i(\vg - \hat\vg_i)$. Show that the Hessian of
    $-\sum_i\ln L_i(\vg_i(\kappa))$ with respect to the pixels is
    $\sum_i J_i^\T F_iJ_i + (\text{terms} \propto \text{residuals} \times \nabla\nabla\vg_i)$, with
    $J_i = \partial\vg_i/\partial\kappa$.
  \item Why are the dropped terms small for us? (Think about how $\vg_i$ depends on $\kappa$; see Part D.)
  \item A bilinear interpolant is linear along each axis inside a grid cell. What is its second derivative along $g_1$
    there? Why does that make the \emph{exact} Hessian of our objective useless, and why does $J^\T FJ$ not have this
    problem?
\end{enumerate}

\solution{
(a) The gradient is $-\sum_i J_i^\T\nabla_{\vg}\ln L_i = \sum_i J_i^\T F_i(\vg_i - \hat\vg_i)$. Differentiating again gives
$\sum_i J_i^\T F_iJ_i$ plus terms with $\partial J_i/\partial\kappa$ multiplied by the residual $F_i(\vg_i - \hat\vg_i)$.
This is the weighted version of MNT04's Eq.~(3.4b).
(b) $\vg_i$ is linear in $\kappa$ except for the $1/(1-\kappa)$ factor (Part D), so $\nabla\nabla\vg_i$ is small, and it is
multiplied by residuals that scatter around zero.
(c) Zero. Inside each grid cell the interpolant has no curvature along $g_1$ or $g_2$; the curvature lives only in the
kinks between cells. The exact Hessian therefore misses most of the data's information and can even have negative
directions (in our first attempt, 12 pixels got $\sigma = 0$). $J^\T FJ$ takes the curvature from each galaxy's known
precision $F_i$, which is positive definite by construction.
}

% --------------------------------------------------------------------
\problem{F3}{Two different uncertainties, and the resolution matrix}

\begin{given}
In the linear-Gaussian model of E2, write $\mathsf H_d = A^\T N^{-1}A$ (data) and $\mathsf H_p = S^{-1}$ (prior), so
$\mathsf H = \mathsf H_d + \mathsf H_p$ and $\hat{\mathbf s} = \mathsf H^{-1}A^\T N^{-1}\vd$. Let the truth be
$\mathbf s_{\rm true}$, with $\vd = A\mathbf s_{\rm true} + \mathbf n$. \emph{Derived here}
(\texttt{KAPPA\_FITTING\_EXPLAINED.md}, Sec.~8.3).
\end{given}

\need{E2, F1.}

\begin{enumerate}
  \item Show that $\E_{\mathbf n}[\hat{\mathbf s}] = \mathsf R\,\mathbf s_{\rm true}$ with $\mathsf R = \mathsf H^{-1}\mathsf H_d$.
  \item Show that the covariance of $\hat{\mathbf s}$ over noise realizations is $\mathsf H^{-1}\mathsf H_d\mathsf H^{-1}$.
  \item Scalar case ($\mathsf H_d = h$, $\mathsf H_p = p$): compare this noise variance with the posterior variance
    $1/(h+p)$. Which is larger, and what does each one answer?
  \item On our map, the noise covariance predicts an aperture scatter of 0.0034 and the posterior covariance 0.0067. In
    the scalar picture, what data-to-prior ratio $h/p$ does that imply? For ``is this blob a noise fluctuation?'',
    which one should the S/N map use?
\end{enumerate}

\solution{
(a) $\E[\vd] = A\mathbf s_{\rm true}$, so \[\E[\hat{\mathbf s}] = \mathsf H^{-1}A^\T N^{-1}A\,\mathbf s_{\rm true} = \mathsf H^{-1}\mathsf H_d\,\mathbf s_{\rm true}.\]
(b) $\mathrm{Cov} = \mathsf H^{-1}A^\T N^{-1}\,N\,N^{-1}A\mathsf H^{-1} = \mathsf H^{-1}\mathsf H_d\mathsf H^{-1}$.
(c) Noise variance $h/(h+p)^2$ = posterior variance $\times\, h/(h+p)$, so it is always smaller. The posterior variance
answers ``how unsure am I about the true value?'', including what the prior allows. The noise variance answers ``how
much would this map change with new noise?''.
(d) $(0.0034/0.0067)^2 = 0.26 = h/(h+p)$, so $h/p \approx 0.35$: on the subhalo's scale the prior outweighs the data
about 3 to 1. The S/N map should use the noise covariance, which is what checks out against the scatter over
synthetic realizations (0.0029).
}

% --------------------------------------------------------------------
\problem{F4}{\py\ How many numbers do the data actually determine?}

\begin{given}
In the eigenbasis of the data Hessian (eigenvalues $\lambda_a$), with an isotropic prior of precision $\alpha$, the number
of parameters well determined by the data is $\gamma = \sum_a \lambda_a/(\lambda_a + \alpha)$. Each term, between 0 and
1, measures the data's strength relative to the prior in that direction.
\source{\cMacKayGamma.}
\end{given}

\need{F3 (the resolution matrix $\mathsf R$), E5 (the 1D toy problem).}

\begin{enumerate}
  \item Show that $\mathrm{tr}\,\mathsf R = \gamma$ when $\mathsf H_p = \alpha I$. (Diagonalise $\mathsf H_d$.)
  \item \py\ For the 1D problem of E5, compute $\mathrm{tr}\,\mathsf R$ as a function of the number of pixels (keeping
    the physical $\ell$ fixed). Does it keep growing?
  \item Our tuned map has 144 pixels and about 150 galaxies, and $\mathrm{tr}\,\mathsf R = 16.5$. With a weaker prior
    ($\sk = 0.1$) it is 53; with 576 pixels it is 10. Answer the question ``how can 141 galaxies fit 144 pixels?''
  \item What sets the map's effective resolution, if not the pixel size?
\end{enumerate}

\solution{
(a) With $\mathsf H_d = U\Lambda U^\T$, $\mathsf R = (\mathsf H_d + \alpha I)^{-1}\mathsf H_d = U\,\mathrm{diag}\big(\lambda_a/(\lambda_a+\alpha)\big)U^\T$,
whose trace is $\gamma$.
(b) No. In the notebook $\mathrm{tr}\,\mathsf R$ rises to $\approx 19$ near 60 pixels and then \emph{falls} (7.7 at
240). Extra pixels add only directions the data cannot see. And with $\sk$ fixed per pixel, the amplitude term
$\sum\kappa_j^2$ runs over more pixels, so every smooth pattern is penalised more: a finer grid with the same $\sk$ is a
stronger prior.
(c) They don't, and they don't need to. It is a regularised fit: the galaxies pin down about 16 independent patterns,
and the prior fixes the remaining $\sim128$. A weaker prior lets the data set more patterns, but noisily. More pixels
add only prior-dominated directions (and, at fixed $\sk$, a stronger prior), as the drop to 10 at 576 pixels shows.
(d) The rows of $\mathsf R$: their width is the resolution and their height the shrinkage, both set jointly by the
galaxies (number, positions, precisions) and by the prior. The pixel size only sets the finest scale the map could
show.
}

\section*{Part G. The cluster analysis}\addcontentsline{toc}{section}{Part G. The cluster analysis}
% ============================== Part G ==============================
\problem{G1}{Bernstein \& Armstrong for a constant shear}

\begin{given}
BA14 write the posterior of one shear shared by many galaxies as $P(\vg\mid D) \propto P(\vg)\prod_iP(D_i\mid\vg)$. They
expand each galaxy's likelihood to second order about $\vg = 0$,
$P(D_i\mid\vg) \approx P_i + \vg\cdot\mathbf Q_i + \tfrac12\vg\cdot\mathbf R_i\cdot\vg$ (here $\mathbf R_i$ is BA14's
second-derivative matrix, not our rotation), and find a Gaussian posterior with
$C_g^{-1} = \sum_i(\mathbf Q_i\mathbf Q_i^\T/P_i^2 - \mathbf R_i/P_i)$ and $\bar\vg = C_g\sum_i\mathbf Q_i/P_i$.
\source{\cBAconst.}
\end{given}

\need{A3, A4, B3, B4.}

\begin{enumerate}
  \item Show that $\mathbf Q_i/P_i = \nabla\ln P(D_i\mid\vg)|_0$ and
    $\mathbf Q_i\mathbf Q_i^\T/P_i^2 - \mathbf R_i/P_i = -\nabla\nabla\ln P(D_i\mid\vg)|_0 \equiv \mathbf A_i$.
  \item For a Gaussian likelihood $\ln L_i = -\tfrac12(\vg - \hat\vg_i)^\T F_i(\vg - \hat\vg_i)$, show that $\mathbf A_i = F_i$
    and $\nabla\ln L_i|_0 = F_i\hat\vg_i$, and hence that BA14's $\bar\vg$ is our stack $\Ftot^{-1}\sum_iF_i\hat\vg_i$ (A3).
  \item Our likelihood is $L_i = q_i/\pitrain$ (A4) on a $401^2$ grid, multiplied without Taylor expansion. What do we gain
    over the quadratic form, and what does the Gaussian form still do for us?
\end{enumerate}

\solution{
(a) $\nabla\ln P = \nabla P/P$, which at 0 is $\mathbf Q_i/P_i$. $\nabla\nabla\ln P = \nabla\nabla P/P - \nabla P\nabla P^\T/P^2$,
which at 0 is $\mathbf R_i/P_i - \mathbf Q_i\mathbf Q_i^\T/P_i^2$.
(b) $\nabla\ln L_i = -F_i(\vg - \hat\vg_i)$, which at 0 is $F_i\hat\vg_i$, and $-\nabla\nabla\ln L_i = F_i$. So
$C_g^{-1} = \sum_iF_i$ and $\bar\vg = (\sum_iF_i)^{-1}\sum_iF_i\hat\vg_i$.
(c) The full product is exact at any $|g| < 0.1$ and includes the higher-order terms BA14 add as perturbations
(\cBAthird). The Gaussian form still defines the reference shear: BA14's $\mathbf A_i$ is our $F_i$, so $\gref$ uses
the same weights the stack does (B4).
}

% --------------------------------------------------------------------
\problem{G2}{Bernstein \& Armstrong for a varying shear}

\begin{given}
For any model with parameters $\mathbf a$ that predicts a shear $\vg_i(\mathbf a)$ at each galaxy, BA14's quadratic
expansion gives (written per galaxy)
\[
-\ln P(\mathbf a\mid D) \approx \text{const} - \ln P(\mathbf a) - \sum_i\vg_i(\mathbf a)\cdot\frac{\mathbf Q_i}{P_i}
+ \frac12\sum_i\vg_i(\mathbf a)\cdot\mathbf A_i\cdot\vg_i(\mathbf a).
\]
If $\vg_i$ is linear in $\mathbf a$ and the prior is Gaussian, the maximum has a closed form.
\source{\cBAvarying.}
\end{given}

\need{G1, B4, E2, E4.}

\begin{enumerate}
  \item \emph{Cells.} Let $\mathbf a$ be one shear per cell, with $\vg_i(\mathbf a) = R_i^\T\vG_{c(i)}$. Show that the
    posterior separates into independent cell stacks. Using $\E[\mathbf Q_i/P_i] \approx \mathbf A_i\vt_i$ (the linear
    response of B4), show that the stack converges to $\gref$.
  \item \emph{$\kappa$ map.} Let $\mathbf a$ be the $\kappa$ pixels, with $\vg_i(\kappa)$ from Part D. Neglecting the
    $1/(1-\kappa)$ factor, $\vg_i$ is linear in $\kappa$: write $\vg = M\kappa$. With the prior of E4, which closed form
    does BA14's statement give, and what is it called?
  \item \emph{Halo fit.} What is $\mathbf a$, and is the model linear?
  \item In one sentence: what does a galaxy contribute to any shear-field model?
\end{enumerate}

\solution{
(a) Each galaxy's term depends only on its own cell's $\vG_c$, so the sum splits by cell, giving one A3/G1 stack per
cell. Setting the gradient with respect to $\vG_c$ to zero, $\sum_{i\in c}\big(R_i\mathbf A_iR_i^\T\vG_c - R_i\mathbf Q_i/P_i\big) = 0$,
so $\hat\vG_c = (\sum_iF_i')^{-1}\sum_iR_i\mathbf Q_i/P_i$ with $F_i' = R_i\mathbf A_iR_i^\T$. With
$\E[\mathbf Q_i/P_i] = \mathbf A_i\vg_i = \mathbf A_iR_i^\T\vt_i$, each term becomes $F_i'\vt_i$, and
$\E[\hat\vG_c] = (\sum_iF_i')^{-1}\sum_iF_i'\vt_i = \gref$.
(b) Write the linear forward model as $\vg = M\kappa$ ($M$ stacks the rows $\partial\vg_i/\partial\kappa$), let
$\mathcal A = \mathrm{blockdiag}(\mathbf A_i)$ and $\mathbf q$ stack the $\mathbf Q_i/P_i$, and take the Gaussian prior
$S^{-1}$ of E4. The objective is quadratic, so $\hat\kappa = (S^{-1} + M^\T\mathcal AM)^{-1}M^\T\mathbf q$. This is E2's
estimator with $\mathcal A$ in the role of $N^{-1}$ (and $\mathbf q = \mathcal A\hat\vg$ for Gaussian likelihoods): the
Wiener filter. Our map is its non-linear, non-Gaussian generalisation, solved numerically.
(c) $\mathbf a = (M_{\rm main}, M_{\rm sub})$; the NFW shear is non-linear in the masses, so it needs sampling or
optimisation.
(d) Its log-likelihood evaluated at the shear the model predicts at that galaxy's position: no per-galaxy ``shear
estimate'' is ever needed.
}

% --------------------------------------------------------------------
\problem{G3}{Why the difference map is clean}

\begin{given}
In each realization, the main-only and main + sub datasets have the same galaxies and the same noise draws; only the
true shears differ. The map estimator is (nearly) the linear MAP of E2, $\hat\kappa = \mathsf H^{-1}A^\T N^{-1}\vd$. The
difference map is $\hat\kappa_{\rm main+sub} - \hat\kappa_{\rm main}$. \emph{Derived here}
(\texttt{KAPPA\_FITTING\_EXPLAINED.md}, Sec.~9.4).
\end{given}

\need{E2, F3, F4.}

\begin{enumerate}
  \item In the linear model, show that the difference map equals $\mathsf R\,\kappa_{\rm subhalo}$ \emph{exactly}, with
    no noise at all.
  \item Our real difference maps have a background rms of $\approx0.002$ against a peak of $\approx0.014$. Where can
    noise come from, if (a) says there is none?
  \item Why can't a real observation make a difference map, and what does it show instead?
  \item The subhalo's true peak pixel is 0.030 and the difference map peaks at about 0.012--0.014. Explain with
    $\mathsf R$.
\end{enumerate}

\solution{
(a) $\vd_{\rm sub} - \vd_{\rm main} = A(\kappa_{\rm sub} - \kappa_{\rm main}) = A\kappa_{\rm subhalo}$, because the noise is
identical and cancels. Then the difference map is $\mathsf H^{-1}A^\T N^{-1}A\kappa_{\rm subhalo} = \mathsf R\,\kappa_{\rm subhalo}$.
(b) From the non-linearity: $g = \gamma/(1-\kappa)$ couples the subhalo to the main halo, and the NPE likelihoods are not
exactly Gaussian, so the noise does not cancel perfectly.
(c) A real cluster gives one map; there is no noise-matched twin without the subhalo. The difference map shows how
visible the subhalo is \emph{through the method}: its position, size and amplitude after the prior's smoothing. The
significance a real map could claim comes from G4.
(d) The rows of $\mathsf R$ near the subhalo are broad and low (F4(d)): the map shrinks and spreads it to about 40\% of
its peak.
}

% --------------------------------------------------------------------
\problem{G4}{\py\ The significance one map can claim}

\begin{given}
For each of $K$ independent realizations of the cluster (new galaxies, positions and noise), we average the
reconstructed $\kappa$ over a 0.5\,Mpc disk at the subhalo position: $a_{{\rm sub},r}$ for the main + sub map,
$a_{{\rm main},r}$ for the main-only map. The single-map significance is
$z_r = (a_{{\rm sub},r} - \overline{a_{\rm main}})/{\rm sd}(a_{\rm main})$, with the mean and sd over realizations. An
aperture average of a mass map is a standard detection statistic (\cSchneiderMap); the empirical-null $z$ is
\emph{derived here}.
\end{given}

\need{B1, G3.}

\begin{enumerate}
  \item Show that subtracting the same fixed map (for example the true main-halo $\kappa$) from every map leaves every
    $z_r$ unchanged.
  \item Why is $z_r$ computed against the other realizations' main-only maps, and not against realization $r$'s own
    main-only map?
  \item Our null is built from maps of the \emph{true} main halo. For a real cluster it would come from simulations of
    a \emph{fitted} main halo. Why does that make our $z$ an upper bound?
  \item With our numbers, $\overline{a_{\rm main}} = 0.0072$, ${\rm sd} = 0.0038$ and $\overline{a_{\rm sub}} = 0.0156$,
    compute the $z$ of the average realization.
  \item \py\ In the 1D toy of E5, generate 50 realizations with and without a narrow bump, compute $z_r$, and estimate
    the fraction with $z > 2$.
\end{enumerate}

\solution{
(a) Subtracting a fixed map subtracts the same constant $c$ (its aperture average) from every $a_{{\rm sub},r}$ and every
$a_{{\rm main},r}$. The difference in the numerator and the sd are unchanged.
(b) Realization $r$'s own main-only map shares its noise, which would cancel (G3), and a real observation has no such
twin. The other realizations represent ``what that spot looks like in a cluster without a subhalo'', across noise and
galaxy layouts.
(c) Uncertainty in the fitted mass, concentration and centre adds scatter to the null, widening its sd and lowering $z$.
(d) $(0.0156 - 0.0072)/0.0038 = 2.2$ (the median over realizations is 2.4).
(e) The notebook (one seed) gives median $z \approx 1.9$, with 48\% of realizations above $z = 2$: a toy version of our 2.4 and 60\%. It also confirms (a): subtracting the truth changes no $z$.
}

% --------------------------------------------------------------------
\problem{G5}{Reading the ensemble numbers}

\begin{given}
Our results over 50 independent NPE realizations (\texttt{repo\_reports/PAPER\_DRAFT\_HANDOFF.md}): one realization
constrains $m$ to $\approx\pm0.3$; the pooled ensemble gives $m = +0.04 \pm 0.04$ (polar $g_t$); the pull spread is 0.84;
the likelihood map's median $z$ is 2.4 (KS: 1.1); and the best possible matched-filter S/N (C5(d), with full precision
matrices) has median 3.7 over realizations. Statistical errors average down as $1/\sqrt K$ (A2).
\end{given}

\need{A2, B1, C5.}

\begin{enumerate}
  \item Check $0.3/\sqrt{50}$ against the ensemble error. How many realizations would $\pm0.01$ need? Why doesn't the
    cluster test aim for that?
  \item What fraction of the matched-filter S/N does the likelihood map achieve, and what fraction does KS achieve?
    Name two reasons the map falls short of the matched filter.
  \item With 4 times the source density (Xu sample 1), roughly what median $z$ would you expect, if $z \propto \sqrt{n}$?
\end{enumerate}

\solution{
(a) $0.3/\sqrt{50} = 0.042$, consistent. $\pm0.01$ needs $(0.3/0.01)^2 = 900$ realizations, because one cluster's cells
span only $\approx0.04$ in shear and the slope is poorly pinned. The test asks a different question: no extra bias
from shear gradients at the few-percent level. Percent-level $m$ comes from the constant-shear stacks.
(b) $2.4/3.7 = 65\%$; KS $1.1/3.7 = 30\%$. The matched filter knows the subhalo's exact profile and position, while the
map assumes neither. The map's prior also shrinks the subhalo and the null includes galaxy-layout variance.
(c) $\approx2.4\times2 = 4.8$.
}

% --------------------------------------------------------------------
\problem{G6}{Why the map settings were chosen on synthetic data}

\begin{given}
We tried 27 likelihood-map settings and 4 KS settings. The rule ``highest contrast with a false-peak rate $\le 0.10$''
was fixed before the NPE maps were looked at, and applied to synthetic data only. For a standard normal, $P(z > 2) =
0.0228$. \emph{Derived here}.
\end{given}

\need{B1, G4.}

\begin{enumerate}
  \item If instead we had tried all 31 settings on the real map and reported the best $z$, and the settings were
    independent, what is the chance that at least one gives $z > 2$ when there is no subhalo?
  \item Real settings are correlated. Does that make the problem better or worse than (a)?
  \item Why does choosing on synthetic data, with the rule fixed in advance, avoid the problem?
\end{enumerate}

\solution{
(a) $1 - (1 - 0.0228)^{31} = 0.51$: a coin flip.
(b) Better, since correlated settings behave like fewer independent tries, but the bias stays positive.
(c) The choice never sees the data it is judged on, so the reported $z$ is a single, pre-registered test.
}


\section*{Sources}\addcontentsline{toc}{section}{Sources}
Every section and equation number cited in the Given boxes was read in the source text (or, for scanned and badly
extracted pages, in rendered page images); \texttt{references.md} quotes the line each citation points to. Equation
numbers follow the arXiv versions where those were used.
\begin{description}[itemsep=2pt,font=\normalfont\bfseries]
\item[BA14] Bernstein, G.~M. \& Armstrong, R. 2014, \emph{Bayesian lensing shear measurement}, MNRAS 438, 1880. arXiv:1304.1843.
\item[B\&S01] Bartelmann, M. \& Schneider, P. 2001, \emph{Weak gravitational lensing}, Phys.\ Rep.\ 340, 291. arXiv:astro-ph/9912508.
\item[HBL10] Hogg, D.~W., Bovy, J. \& Lang, D. 2010, \emph{Data analysis recipes: Fitting a model to data}. arXiv:1008.4686.
\item[H12] Hogg, D.~W. 2012, \emph{Data analysis recipes: Probability calculus for inference}. arXiv:1205.4446.
\item[HMB10] Hogg, D.~W., Myers, A.~D. \& Bovy, J. 2010, \emph{Inferring the eccentricity distribution}, ApJ 725, 2166. arXiv:1008.4146.
\item[ITILA03] MacKay, D.~J.~C. 2003, \emph{Information Theory, Inference, and Learning Algorithms}, Cambridge University Press (free PDF from the author).
\item[J18] Jeffrey, N. et al. 2018, \emph{Improving weak lensing mass map reconstructions using Gaussian and sparsity priors: application to DES SV}, MNRAS 479, 2871. arXiv:1801.08945.
\item[Kilbinger15] Kilbinger, M. 2015, \emph{Cosmology with cosmic shear observations: a review}, Rep.\ Prog.\ Phys.\ 78, 086901. arXiv:1411.0115.
\item[KS93] Kaiser, N. \& Squires, G. 1993, \emph{Mapping the dark matter with weak gravitational lensing}, ApJ 404, 441. (Read from the scanned ADS copy. KS93's notation differs: their $\Sigma = 2\kappa$, their potential has the opposite sign, and their wavevector is called $\kappa$.)
\item[KvdV12] Kleijn, B.~J.~K. \& van der Vaart, A.~W. 2012, \emph{The Bernstein--von-Mises theorem under misspecification}, Electron.\ J.\ Statist.\ 6, 354.
\item[M92] MacKay, D.~J.~C. 1992, \emph{Bayesian interpolation}, Neural Computation 4, 415 (equation number from the author's preprint).
\item[M92t] MacKay, D.~J.~C. 1992, \emph{Bayesian Methods for Adaptive Models}, PhD thesis, California Institute of Technology.
\item[MNT04] Madsen, K., Nielsen, H.~B. \& Tingleff, O. 2004, \emph{Methods for Non-Linear Least Squares Problems}, 2nd ed., IMM, Technical University of Denmark.
\item[S96] Schneider, P. 1996, \emph{Detection of (dark) matter concentrations via weak gravitational lensing}, MNRAS 283, 837. arXiv:astro-ph/9601039.
\item[Sch15] Schneider, M.~D., Hogg, D.~W., Marshall, P.~J. et al. 2015, \emph{Hierarchical probabilistic inference of cosmic shear}, ApJ 807, 87. arXiv:1411.2608.
\item[WB00] Wright, C.~O. \& Brainerd, T.~G. 2000, \emph{Gravitational lensing by NFW halos}, ApJ 534, 34. arXiv:astro-ph/9908213.
\item[Z95] Zaroubi, S., Hoffman, Y., Fisher, K.~B. \& Lahav, O. 1995, \emph{Wiener reconstruction of the large-scale structure}, ApJ 449, 446. arXiv:astro-ph/9410080.
\end{description}
Not checked: White, H. 1982, \emph{Maximum likelihood estimation of misspecified models}, Econometrica 50, 1 (not openly available); KvdV12 is cited in its place.

